\chapter{Integrite scientifique, IA generative et reproductibilite}

\section{Declaration d'usage de l'IA generative}
\guided{L'usage d'un outil d'IA pour reformuler, coder, g\'en\'erer des figures, proposer des plans ou synth\'etiser des ressources doit \^etre explicitement d\'eclar\'e. L'\'etudiant demeure enti\`erement responsable de l'exactitude du contenu et du respect des r\`egles de citation.}

\begin{quote}
Exemple de formulation : \textit{Des outils d'IA g\'en\'erative ont \`a certains moments assist\'e la reformulation de passages, la suggestion d'une structure de document ou la g\'en\'eration d'exemples de code. L'ensemble des contenus a \'et\'e v\'erifi\'e, corrig\'e et valid\'e par l'auteur, qui en assume l'enti\`ere responsabilit\'e scientifique et acad\'emique.}
\end{quote}

\section{Reproductibilite minimale attendue}
\begin{itemize}[leftmargin=1.2cm]
\item Jeux de donn\'ees : source, taille, date, mode d'acquisition, restrictions.
\item Outils : versions logicielles, biblioth\`eques, environnement mat\'eriel, configuration syst\`eme.
\item Param\`etres critiques : hyperparam\`etres, seuils, fr\'equences, dimensions, protocoles.
\item El\'ements de validation : scripts, cahiers de tests, captures de configuration, sch\'emas, nomenclatures.
\item Confidentialit\'e : distinguer ce qui est diffusable, masqu\'e ou anonymis\'e.
\end{itemize}

\section{Rappels de conformite}
\begin{itemize}[leftmargin=1.2cm]
\item \okitem{} Toute figure, tableau, code ou id\'ee emprunt\'e doit \^etre attribu\'e.
\item \okitem{} Les r\'esultats doivent \^etre interpr\'et\'es avec esprit critique.
\item \warnitem{} Un taux de similarit\'e trop \^elev\'e compromet la recevabilit\'e du rapport.
\item \warnitem{} Un texte fluent n'est pas une preuve scientifique s'il n'est pas reli\'e \`a une m\'ethode et \`a des r\'esultats v\'erifiables.
\end{itemize}
